\section{คู่มือสำหรับอาจารย์ประจำวิชาและเจ้าหน้าที่}
\subsection{อาจารย์ประจำวิชา}
ผู้มี \role{LECTURER} เปิดเมนู \textbf{อาจารย์ประจำวิชา} badge แสดงคำร้องอนุมัติแล้วที่ยังไม่เปิดอ่านใน section ที่สอน เปิดคำร้องเพื่อตรวจนิสิต วันเวลา เหตุผล และข้อมูลประกอบ ระบบบันทึกการเปิดอ่าน บทบาทนี้เป็นการรับทราบ ไม่ใช่ขั้นอนุมัติหลัก เมนู \textbf{วิชาที่สอน} ใช้ตรวจ section ที่นำมาจับคู่

\subsection{เจ้าหน้าที่: รายการใบลา}
ผู้มี \role{ADMIN} เปิด \textbf{รายการใบลาภาควิชา} ใช้ badge และตัวกรอง ตรวจสถานะ เส้นทางอนุมัติ และข้อมูลผิดปกติ หลีกเลี่ยงการแก้ master data ระหว่างวิเคราะห์โดยไม่บันทึกเหตุผล

\subsection{เจ้าหน้าที่: ภาคการศึกษา}
เปิด \textbf{จัดการปีการศึกษา} เพิ่มหรือแก้ภาคและช่วงวันที่ แล้วกำหนดภาคปัจจุบัน หลังเปลี่ยนควรทดสอบคำร้องในสภาพแวดล้อมที่เหมาะสมและตรวจ snapshot

\subsection{เจ้าหน้าที่: คณะและสาขา}
เปิด \textbf{จัดการคณะ/สาขา} ตรวจ code และชื่อภาษาอังกฤษให้ตรงระบบต้นทาง เพราะใช้หาเส้นทางอนุมัติ ก่อนลบต้องตรวจว่าไม่มีนิสิตหรือข้อมูลอื่นอ้างอิง

\subsection{เจ้าหน้าที่: ตั้งค่าระบบ}
หน้า \textbf{ตั้งค่าระบบ} ใช้กำหนดกฎอนุมัติอัตโนมัติแยกตาม role ต้องได้รับอนุมัติตามนโยบายและตรวจผลหลังแก้ ไม่ควรเปิดเพื่อทดลองใน production

\subsection{รายงาน}
รายงานขึ้นกับ role และข้อมูล warehouse สามารถกรองตามปี ภาค สถานะ ประเภทลา หน่วยงาน รายวิชา หรือนิสิต หากตัวเลขไม่ตรงรายการล่าสุดให้ตรวจรอบนำเข้าข้อมูล
